\begin{textsample}{NEG Dim 4 – 2013\_c6 – Score 8.00 – 2013\_c6/enade\_2013\_serv\_soc\_b0ca85ce74.txt}  \label{ex:f4_neg_316}
CAPÍTULO 3 ANÁLISE TÉCNICA DA PROVA Este capítulo tem por objetivo apresentar o desempenho dos estudantes concluintes de Serviço Social no ENADE/2013.
Para isso, foram calculadas as estatísticas básicas da prova em seu todo ( seção 3.1.1 ), bem como as estatísticas dos componentes relacionadas ao Componente de Formação Geral ( seção 3.1.2 ) e ao de Conhecimento Específico da Área ( seção 3.1.3 ).
Dadas as suas características, foram analisadas em separado, as questões objetivas ( seção 3.2 ) e as questões discursivas ( seção 3.3 ).
Para as questões objetivas, foram disponibilizados os índices de facilidade e o de discriminação bisserial, também em separado para os Componentes de Formação Geral ( seção 3.2.1 ) e de Conhecimento Específico ( seção 3.2.2 ).
De cada componente, uma das questões foi escolhida para exemplificar a análise gráfica relacionando as alternativas escolhidas pelos estudantes ( inclusive o gabarito ) com o número de acertos no componente.
O Anexo I apresenta a íntegra da análise gráfica para todas as questões objetivas.
Para cada uma das questões discursivas, os conteúdos dos tipos mais comuns de respostas dos alunos são apresentados e comparados com o padrão de resposta esperado ( ver Anexo VIII com o padrão de respostas ).
Tomando como base as duas questões discursivas do Componente de Formação Geral, a seção 3.3.1.6 apresenta comentários sobre a correção das respostas com respeito à Língua Portuguesa.
Nas tabelas, são apresentados o tamanho da população inscrita e de presentes, e as seguintes estatísticas das notas11 : média do desempenho na prova, erro padrão da média, desvio padrão, nota mínima, mediana e nota máxima.
Tais estatísticas contemplam o total de estudantes concluintes da Área de Serviço Social inscritos e presentes à prova do ENADE/2013, tendo -se em vista as seguintes agregações : ( a ) as Grandes Regiões e o país como um todo ; ( b ) a Categoria Administrativa ; e ( c ) a Organização Acadêmica.
Em relação aos gráficos de distribuição de notas, o intervalo considerado foi de 10 unidades, aberto à esquerda e fechado à direita, com exceção do primeiro intervalo, [ 0 ; 10 ], fechado em ambos os extremos.
Para os gráficos de distribuição das notas das questões discursivas, foram consideradas mais duas categorias : questão em branco e nota zero. 11 Essas estatísticas e outras estão definidas no Capítulo 1. 3.1 ESTATÍSTICAS BÁSICAS DA PROVA 3.1.1 Estatísticas Básicas Gerais A Tabela 3.1 apresenta as estatísticas básicas da prova por grande Região.
A população total de inscritos foi de 41.716.
Destes, 30.857 estiveram presentes, sendo 26,0 \% o índice de não comparecimento.
A Região de maior abstenção foi a Norte ( 75,3 \% ) e a de menor abstenção foi a Nordeste ( 12,0 \% ).
A média das notas da prova como um todo ( nas seções seguintes serão analisados os Componentes de Formação Geral e de Conhecimento Específico ) foi 36,4, sendo que os alunos da região Centro-Oeste obtiveram a média mais baixa ( 33,3 ) e os da região Nordeste obtiveram a média mais alta ( 40,8 ).
As demais médias foram : 37,3 na região Norte, 37,9 na região Sudeste e 34,5 na região Sul.
O desvio padrão para o Brasil como um todo foi 14,6, sendo o maior desvio padrão encontrado na região Nordeste ( 16,7 ) e o menor na região Centro-Oeste ( 12,1 ), indicando uma menor dispersão das notas desta última região.
A região que obteve a maior nota máxima foi a Nordeste ( 89,3 ), ao passo que a região que atingiu a menor nota máxima foi a Norte ( 82,2 ).
A mediana do Brasil como um todo foi 34,8, sendo a maior mediana obtida na região Nordeste ( 40,6 ) e a menor obtida na Centro-Oeste ( 31,8 ).
A nota mínima foi zero em todas as regiões, sem exceção.
O comportamento das notas dos estudantes de todo o Brasil pode ser observado no Gráfico 3.1 que apresenta um histograma com a distribuição das mesmas.
Essa é uma distribuição unimodal com moda no intervalo ( 30 ; 40 ].
Destaca -se, também, o intervalo ( 20 ; 30 ] com frequência relativa próxima à modal.
O coeficiente de assimetria da distribuição das notas é positivo ( 0,33 ) e não podemos considerar que a distribuição é aproximadamente simétrica.
As distribuições por Grande Região também apresentam assimetria positiva, concentração pouco maior do lado esquerdo do histograma e mais espalhada do lado direito, excetuando -se a região Nordeste que possui assimetria negativa.
Os Gráficos 3.2, 3.3 e 3.4 apresentam informações referentes à média da nota final dos participantes, desagregando os resultados de acordo com, respectivamente, as Grandes Regiões do país, a Categoria Administrativa e a Organização Acadêmica.
Os gráficos apresentam o valor da média das notas como uma barra e os extremos do intervalo de confiança de 95 \% como linhas verticais unidas por uma linha horizontal na forma da letra H maiúscula.
Considerando -se o gráfico de notas segundo Grande Região ( Gráfico 3.2 ), observa -se que existe diferença estatisticamente significativa ao nível de 95 \% entre a maior média, obtida na região Nordeste ( 40,8 ) e a menor, obtida na região Centro-Oeste ( 33,3 ).
Levando -se em conta os estudantes por Categorias Administrativas ( Gráfico 3.3 ), observa -se que existe diferença estatisticamente significante entre as médias das notas das IES Públicas e Privadas.
Pode ser observado ainda que a média dos alunos de IES Públicas ( 40,6 ) é maior que a dos alunos de IES Privadas ( 35,8 ).
A diferença entre as médias das regiões Nordeste e Centro-Oeste ( 7,5 ), a maior e a menor média, é superior a diferença entre IES Públicas e Privadas ( 4,8 ).
Tendo como foco o Gráfico 3.4, que apresenta as notas médias das provas segundo Organização Acadêmica, constata -se que a diferença da média das notas dos estudantes provenientes de Universidades é estatisticamente significativa ao nível de 95 \% das médias das notas dos estudantes provenientes de Centros Universitários e Faculdades, assim como a diferença entre a média dos estudantes de Centros Universitários e Faculdades.
A menor média foi obtida pelos estudantes de Universidades ( 35,4 ), e a maior, pelos de Faculdades ( 39,6 ). 3.1.2 Estatísticas Básicas no Componente de Formação Geral A Tabela 3.2 apresenta as estatísticas básicas em relação ao componente da prova que avalia a Formação Geral dos estudantes concluintes.
Os alunos de todo Brasil obtiveram desempenho médio de 40,9.
Quanto à variabilidade, o desvio padrão das notas dos estudantes do Brasil como um todo foi 15,5.
A maior média foi obtida na região Sul ( 41,7 ), e a menor, na região Sudeste ( 39,8 ).
As demais médias foram : 41,2 na região Norte, 40,8 na região Nordeste e 41,0 na região Centro-Oeste.
Já o maior desvio padrão foi obtido na região Nordeste ( 17,1 ) e o menor na região Norte ( 14,2 ).
Os demais desvios padrões foram : 16,9 na região Sudeste, 14,4 na região Sul e 14,3 na região Centro-Oeste.
A maior nota no Componente de Formação Geral da prova do ENADE foi obtida por pelo menos um aluno da região Sul ( 93,6 ) enquanto que a menor nota máxima foi obtida na região Norte ( 82,6 ).
Nas outras regiões as notas máximas foram : 91,0 para a região Nordeste, 92,4 para a região Sudeste e 88,8 para a região Centro-Oeste.
A mediana do Brasil como um todo foi 41,4, sendo a menor mediana encontrada na região Sudeste ( 40,3 ) e a maior encontrada na região Nordeste ( 41,9 ).
A nota mínima nesta parte foi zero em todas as regiões, sem exceção.
O Gráfico 3.5 propicia a avaliação do desempenho dos estudantes no Componente de Formação Geral a partir do histograma da distribuição das notas correspondentes.
A distribuição é unimodal, com moda em ( 40 ; 50 ], enquanto na prova como um todo a moda foi alcançada no intervalo ( 30 ; 40 ].
Nota -se, ainda, que no Gráfico 3.5 as notas apresentam uma maior dispersão do que no Gráfico 3.1 ( distribuição das notas da prova ), confirmado pela comparação dos desvios padrões : 14,6 para a nota da prova como um todo e 15,5 para o Componente de Formação Geral.
Para o Componente de Formação Geral, o coeficiente de assimetria da distribuição das notas dos estudantes é negativo ( -0,19 ).
Por ser um coeficiente pequeno, a distribuição apresenta uma leve concentração à direita e cauda maior à esquerda, porém podemos considerar que a distribuição é aproximadamente simétrica.
Em quase todas as Grandes Regiões os histogramas também possuem assimetria negativa ( entre -0,30 para o Nordeste e -0,06 para o Centro-Oeste ), excetuando -se na região Norte ( 0,11 ) que possui assimetria positiva.
Nos Gráficos 3.6, 3.7 e 3.8 são apresentadas as informações referentes ao desempenho dos concluintes no Componente de Formação Geral, em diferentes agregações : Grande Região do país, Categoria Administrativa e Organização Acadêmica.
Observa -se pelo Gráfico 3.6 que existe diferença estatisticamente significativa ao nível de 95 \% entre a maior média das notas na Componente de Formação Geral, obtida na região Sul ( 41,7 ) e a menor, obtida na região Sudeste ( 39,8 ).
Vemos que o intervalo de confiança mais largo é o da região Norte ; já o intervalo mais estreito é observado na região Sul.
Este fato está relacionado, também, com o tamanho da população presente envolvida, menor na região Norte do que na Sul.
No Gráfico 3.7, que representa as notas médias no Componente de Formação Geral segundo Categoria Administrativa do país, observa -se que existe diferença estatisticamente significativa entre as médias.
Diferente de como ocorreu para as notas da prova como um todo, para o Componente de Formação Geral, os concluintes de Serviço Social, das IES Públicas ( 36,2 ) obtiveram uma média menor do que os das IES Privadas ( 41,6 ).
Considerando -se o tipo de Organização Acadêmica, nota -se, no Gráfico 3.8, que existem diferenças estatisticamente significativas entre os tipos de Organização Acadêmica.
Nas Universidades, a média ( 40,7 ) é menor do que em Centros Universitários e Faculdades, que possuem a mesma média ( 41,4 ). 3.1.3 Estatísticas Básicas do Componente de Conhecimento Específico A Tabela 3.3 apresenta as estatísticas básicas referentes ao Componente de Conhecimento Específico da área de Serviço Social.
A média do desempenho dos alunos do Brasil como um todo foi 34,8.
A maior média foi obtida na região Nordeste ( 40,8 ), e a menor, na região Centro-Oeste ( 30,7 ).
As demais médias foram : 36,1 na região Norte, 37,3 na região Sudeste e 32,1 na região Sul.
Quanto à variabilidade das notas, o desvio padrão do Brasil como um todo foi 16,7, sendo o maior desvio padrão observado na região Nordeste ( 18,8 ) e o menor na região Centro-Oeste ( 13,9 ).
Os demais desvios padrões foram : 16,6 na região Norte, 17,9 na região Sudeste e 14,6 na região Sul.
A mediana das notas dos estudantes de todo o Brasil foi 32,8.
A maior mediana ocorreu na região Nordeste ( 40,0 ) e a menor nas regiões Sul e Centro-Oeste ( 29,8 ).
As demais medianas foram 34,0 na região Norte e 35,8 na região Sudeste.
A nota máxima do Brasil como um todo foi 93,8, sendo obtida por pelo menos um aluno da região Nordeste.
As demais notas máximas foram : 90,5 na região Norte, 89,8 na região Sudeste, 87,8 na região Sul, 90,8 na região Centro-Oeste.
A nota mínima foi zero em todas as regiões, sem exceção.
Assim como os Gráficos 3.1 e 3.5, o Gráfico 3.9, apresentado a seguir, proporciona uma avaliação do desempenho de concluintes em relação ao Componente de Conhecimento Específico com um histograma da distribuição das notas correspondentes.
Esta também é uma distribuição unimodal, e o grupo modal é o ( 20 ; 30 ], abaixo do grupo modal para a Formação Geral.
O coeficiente de assimetria da distribuição das notas do Componente de Conhecimento Específico é positivo ( 0,49 ).
Nota -se pelo histograma ( Gráfico 3.9 ) que esta distribuição apresenta uma leve concentração à esquerda e cauda maior à direita.
Em todas as Grandes Regiões os histogramas também possuem assimetria positiva ( entre 0,07 para o Nordeste e 0,72 para o Centro-Oeste ).
Os Gráficos 3.10, 3.11 e 3.12 apresentam, respectivamente, uma comparação dos resultados em relação à Grande Região do país, à Categoria Administrativa e à Organização Acadêmica, agora levando em conta o desempenho de alunos no Componente de Conhecimento Específico da prova.
Pelo Gráfico 3.10, observa -se que existe diferença estatisticamente significativa entre as médias das notas no Componente de Conhecimento Específico, das regiões Nordeste e Centro-Oeste em relação às demais regiões.
A região Nordeste foi a que obteve média mais elevada ( 40,8 ), enquanto a região Centro-Oeste foi a que obteve média mais baixa ( 30,7 ).
Quanto à Categoria Administrativa ( Gráfico 3.11 ), observa -se um comportamento semelhante àquele da Parte de Formação Geral e à prova como um todo, ou seja, existe diferença estatisticamente significativa entre as médias das IES Públicas e Privadas, sendo que a maior média foi obtida por alunos de IES Públicas de ensino ( 42,0 ).
Quanto ao Gráfico 3.12, observa -se, mais uma vez, que existe diferença estatisticamente significativa ao nível de 95 \% entre as notas no Componente de Conhecimento Específico de Universidades e dos demais tipos de Organização Acadêmica.
Sendo que a média dos concluintes das Universidades ( 33,6 ) foi menor do que de Centros Universitários ( 36,7 ) e de Faculdades ( 39,0 ). 3.2 ANÁLISE DAS QUESTÕES OBJETIVAS 3.2.1 Componente de Formação Geral A Tabela 3.4 apresenta as estatísticas básicas relativas às oito questões objetivas do componente da prova que abrange a Formação Geral dos estudantes.
A média do Brasil foi 43,9.
A menor média foi encontrada na região Norte ( 43,7 ) e a maior na região Nordeste ( 44,3 ).
As demais médias foram 43,8 na região Sudeste e Centro-Oeste e 43,9 na região Sul.
O desvio padrão do Brasil foi 18,3, sendo o maior desvio padrão encontrado na região Nordeste ( 19,6 ) e o menor na região Norte ( 16,7 ).
Os demais desvios foram : 19,4 na região Sudeste, 17,5 na região Sul e 17,3 na região Centro-Oeste.
A mediana do Brasil com um todo foi 50,0, sendo também encontrada nas regiões Nordeste e Sudeste.
Nas regiões Norte, Sul e Centro-Oeste, a mediana foi 37,5.
As notas máximas ( 100,0 ) e notas mínimas ( 0,0 ), nas questões objetivas do Componente de Formação Geral, foram obtidas em todas as regiões, sem exceção.
A Tabela 3.5 apresenta o índice de facilidade e o índice de discriminação ( ponto bisserial ) para cada uma das questões objetivas do Componente de Formação Geral.
Quanto ao índice de facilidade, foram usadas as seguintes cores para diferenciar o nível de dificuldade da questão :  Azul para as questões classificadas com índice muito fácil ( índice \&gt ; =0,86 ), verde para as questões classificadas com índice fácil ( 0,61 a 0,85 ), amarelo para as questões classificadas com médio ( 0,41 a 0,60 ), vermelho para as questões classificadas com difícil ( 0,16 a 0,40 ) e roxo para as questões classificadas com muito difícil ( \&lt ; =0,15 ).
Já quanto ao índice de discriminação, foram usadas as seguintes cores para qualificar a questão :  As questões classificadas com índice fraco receberam a cor vermelha ( índice \&lt ; =0,19 ), as classificadas com médio receberam a cor amarelo ( 0,20 a 0,29 ), as classificadas com bom receberam a cor verde ( 0,30 a 0,39 ) e as classificadas com muito bom ( \&gt ; =0,40 ) receberam a cor azul.
As questões objetivas do Componente de Formação Geral, segundo o índice de facilidade, foram assim avaliadas : das oito questões, uma teve o índice de facilidade classificado como muito fácil e uma como fácil.
Duas questões foram tidas como médio, por terem índice de acertos situado na faixa entre 0,41 e 0,60 ( de 41,0 \% a 60,0 \% de acertos ).
Quatro questões foram consideradas de dificuldade difícil, situando -se no intervalo entre 0,16 e 0,40 do índice de facilidade, ou seja, houve entre 16,0 \% e 40,0 \% de acertos.
Como já comentado, para análise das questões objetivas relativas à Formação Geral segundo o poder de discriminação, utilizou -se, o índice de discriminação ponto bisserial.
Nesta análise, as questões foram assim avaliadas : cinco das oito questões apresentaram A Tabela 3.5 apresenta o índice de facilidade e o índice de discriminação ( ponto bisserial ) para cada uma das questões objetivas do Componente de Formação Geral.
Quanto ao índice de facilidade, foram usadas as seguintes cores para diferenciar o nível de dificuldade da questão :  Azul para as questões classificadas com índice muito fácil ( índice \&gt ; =0,86 ), verde para as questões classificadas com índice fácil ( 0,61 a 0,85 ), amarelo para as questões classificadas com médio ( 0,41 a 0,60 ), vermelho para as questões classificadas com difícil ( 0,16 a 0,40 ) e roxo para as questões classificadas com muito difícil ( \&lt ; =0,15 ).
Já quanto ao índice de discriminação, foram usadas as seguintes cores para qualificar a questão :  As questões classificadas com índice fraco receberam a cor vermelha ( índice \&lt ; =0,19 ), as classificadas com médio receberam a cor amarelo ( 0,20 a 0,29 ), as classificadas com bom receberam a cor verde ( 0,30 a 0,39 ) e as classificadas com muito bom ( \&gt ; =0,40 ) receberam a cor azul.
As questões objetivas do Componente de Formação Geral, segundo o índice de facilidade, foram assim avaliadas : das oito questões, uma teve o índice de facilidade classificado como muito fácil e uma como fácil.
Duas questões foram tidas como médio, por terem índice de acertos situado na faixa entre 0,41 e 0,60 ( de 41,0 \% a 60,0 \% de acertos ).
Quatro questões foram consideradas de dificuldade difícil, situando -se no intervalo entre 0,16 e 0,40 do índice de facilidade, ou seja, houve entre 16,0 \% e 40,0 \% de acertos.
Como já comentado, para análise das questões objetivas relativas à Formação Geral segundo o poder de discriminação, utilizou -se, o índice de discriminação ponto bisserial.
Nesta análise, as questões foram assim avaliadas : cinco das oito questões apresentaram O Gráfico 3.13, para exemplificar, analisa o comportamento da questão de número 3 de Formação Geral.
Trata -se da terceira questão mais difícil e a que obteve o terceiro maior índice de discriminação dessa parte da prova.
Neste gráfico, cada uma das cinco curvas representa o percentual de respostas em determinada alternativa da questão, em função da nota dos estudantes nesta parte da prova ( Formação Geral/Múltipla Escolha ), antes de possíveis eliminações pelo critério do ponto bisserial.
A curva em vermelho corresponde à alternativa E, a correta para esta questão.
Assim, observa -se que entre os estudantes com notas mais baixas, nessa parte do exame, a situação mais frequente foi a escolha de uma das alternativas incorretas : a alternativa C índices acima de 0,40 e, assim, foram classificadas com índice muito bom para esse grupo de alunos ; duas questões tiveram índice de discriminação bom, entre 0,30 e 0,39 ; e uma questão teve o índice de discriminação médio, entre 0,20 e 0,29.
Nenhuma questão foi classificada com um índice fraco e, portanto, nenhuma questão foi eliminada pelo critério ponto bisserial.
O índice de facilidade variou de 0,16 a 0,90, e o de discriminação, de 0,27 a 0,48.
As questões com índice de discriminação muito bom, de números 4, 6 e 8, figuraram entre as mais fáceis desse conjunto : uma ( questão 4 ) classificada na categoria fácil e outras duas na categoria médio do índice de facilidade.
Em particular, a questão 8 apresentou maior poder discriminatório, com índice 0,48, e foi a terceira mais fácil, com uma proporção de 0,50 de acertos.
A questão de número 5 apresentou índice de facilidade 0,16, ou seja, um quantitativo de 16,0 \% de estudantes conseguiu resolvê -la, dentro do universo de participantes.
Além disso, seu índice de discriminação foi médio, com valor igual a 0,27.
O Gráfico 3.13, para exemplificar, analisa o comportamento da questão de número 3 de Formação Geral.
Trata -se da terceira questão mais difícil e a que obteve o terceiro maior índice de discriminação dessa parte da prova.
Neste gráfico, cada uma das cinco curvas representa o percentual de respostas em determinada alternativa da questão, em função da nota dos estudantes nesta parte da prova ( Formação Geral/Múltipla Escolha ), antes de possíveis eliminações pelo critério do ponto bisserial.
A curva em vermelho corresponde à alternativa E, a correta para esta questão.
Assim, observa -se que entre os estudantes com notas mais baixas, nessa parte do exame, a situação mais frequente foi a escolha de uma das alternativas incorretas : a alternativa C ( em preto ), D ( em roxo ) ou A ( azul ).
Na medida em que a nota aumenta, indicando desempenho melhor nesta parte da prova, aumenta concomitantemente a proporção de estudantes que selecionaram a alternativa correta E, atingindo 100 \% para os alunos com 8 acertos.
Essa análise permite verificar como a questão discriminou os grupos de desempenho, justificando o alto índice obtido na questão.
Os gráficos relativos às demais questões de Formação Geral constam do Anexo I. 3.2.2 Componente de Conhecimento Específico A Tabela 3.6 apresenta as estatísticas básicas em relação às questões objetivas do Componente de Conhecimento Específico da prova, por Grande Região.
A média do Brasil deste componente foi de 39,1.
A menor média foi observada na região Centro-Oeste ( 34,6 ) e a maior na região Nordeste ( 45,7 ).
O desvio padrão de todo o Brasil foi 18,6, sendo o menor desvio padrão encontrado na região Centro-Oeste ( 15,5 ) e o maior na região Nordeste ( 21,0 ).
A mediana de todo o Brasil foi 35,0.
Nas regiões Norte, Sul e Centro-Oeste, a mediana foi também 35,0.
Na região Nordeste e na Sudeste, as medianas foram 45,0 e 40,0, respectivamente.
A nota máxima da prova ( 100,0 ) foi obtida, nas questões objetivas do Componente de Conhecimento Específico, por pelo menos um aluno em quase todas as regiões do Brasil, exceto nas regiões Sul e Centro-Oeste ( 95,0 ).
Em todas as regiões, sem exceção, a nota mínima foi zero.
A Tabela 3.7 apresenta o índice de facilidade e o índice de discriminação ( ponto bisserial ) para cada uma das questões objetivas do Componente de Conhecimento Específico.
Para facilitar a diferenciação das questões usaram -se as mesmas cores da Tabela 3.5 para as diferentes classificações dos índices de facilidade e de discriminação.
Dentre as questões objetivas da parte da prova relativa ao Componente de Conhecimento Específico, apenas a questão 35 foi anulada pela Comissão.
Desse modo, a classificação quanto ao índice de facilidade foi estabelecida com base nas 26 questões restantes.
A partir dos índices obtidos, pode -se concluir que a maioria das 26 questões objetivas da prova foi considerada como difícil, 15 ao todo.
As 11 questões restantes foram classificadas quanto ao índice de facilidade da seguinte maneira : três como fácil, quatro como médio e outras quatro como muito difícil.
Já quanto aos índices de discriminação das questões objetivas do Componente de Conhecimento Específico da prova, tem -se como resultado a seguinte classificação : dez das 26 questões foram consideradas como bom, enquanto oito delas tiveram índice de discriminação muito bom.
Assim, para 18 das 26 questões, os índices de discriminação foram bom ou muito bom.
Dentre as demais, duas delas foram classificadas como médio e outras seis como fraco, sendo oito, por conseguinte, a quantidade de questões nos dois patamares mais baixos de discriminação.
Constata -se, assim, que a prova – no que se refere ao Componente de Conhecimento Específico – possuía capacidade mediana de discriminar entre aqueles que dominam ou não o conteúdo.
Dentre as questões que alcançaram os maiores índices de discriminação, as de números 9, 12, 15, 19, 26, 30, 31 e 32 classificadas com índice muito bom, igual ou superior a 0,40, duas delas ( questões 15 e 31 ) foram classificada na categoria fácil, uma delas ( questão 30 ) como médio e outras cinco como difícil, quanto ao índice de facilidade.
A questão de número 10 foi a mais difícil dentre as 26 questões específicas consideradas, com baixo índice de facilidade, apenas 4,0 \% de acertos.
Essa questão apresentou poder discriminatório igualmente baixo, 0,11, o que comprova ter sido esta a mais difícil para os estudantes.
Destaca -se, também, a questão 20, com índice de facilidade 0,08, o que, em termos percentuais, corresponde a 8,0 \% de estudantes que responderam acertadamente, obtendo, ainda, 0,15 de índice de discriminação.
Tais questões foram, portanto, duas das mais difíceis da prova.
As questões 10 e 20 foram eliminadas do cômputo da nota final pelo critério do ponto bisserial.
Além destas duas, as demais questões com índice fraco de discriminação, questões 17, 18, 27 e 33 também não foram computadas.
A título de exemplo das análises do comportamento das questões objetivas, o Gráfico 3.14 analisa a questão 30 do Componente de Conhecimento Específico.
Esta foi uma das questões classificadas como médio nesta prova, apresentando índice de facilidade 0,51, ou seja, 51,0 \% dos estudantes assinalaram acertadamente a opção D, correspondente ao gabarito.
Seu índice de discriminação foi igual a 0,52, classificado como muito bom.
Neste gráfico, cada uma das cinco curvas representa o percentual de respostas em determinada alternativa da questão 30, em função da nota dos estudantes nesta parte da prova, antes de possíveis eliminações de questões pelo critério do ponto bisserial.
A alternativa correta D, representada no gráfico pela curva em roxo, foi escolhida em maiores proporções pelos alunos com desempenho melhor nesta parte da prova.
Já as alternativas incorretas, também denominadas distratores, foram selecionadas principalmente por aqueles com notas mais baixas.
A proporção de alunos que selecionou a resposta correta D aumenta gradativamente, chegando a atingir 100 \% para aqueles que acertaram pelo menos 21 questões, enquanto a proporção dos que escolheram alternativas incorretas decai, a partir da primeira nota não nula, como função da nota nesta parte da prova.
Os gráficos relativos às demais questões do Componente de Conhecimento Específico constam do Anexo I. 3.3 ANÁLISE DAS QUESTÕES DISCURSIVAS 3.3.1 Componente de Formação Geral As análises dos resultados de desempenho dos estudantes de Serviço Social nas duas questões discursivas relativas à Formação Geral encontram -se na Tabela 3.8 e no Gráfico 3.15.
Na Tabela 3.8 observa -se que as notas médias foram mais baixas nesse conjunto de questões do que no das objetivas.
Os estudantes de todo o Brasil obtiveram, em Formação Geral, média 43,9 nas questões objetivas e 36,4 nas questões discursivas.
A mediana também confirma o pior desempenho dos alunos de todo o Brasil nas questões discursivas do Componente de Formação Geral.
Enquanto essa estatística foi de 41,0 para questões discursivas, para as questões objetivas a mediana foi de 50,0.
Pode -se, também, notar um aumento do desvio padrão de 18,3, nas questões objetivas do Componente de Formação Geral dos alunos de todo o Brasil, para 22,3 nas questões discursivas do mesmo componente.
Como já comentado, a mediana de todo o Brasil, neste componente, foi 41,0, sendo esta a maior mediana encontrada na região Sul ( 42,0 ) e a menor na região Sudeste ( 38,5 ).
A maior nota máxima foi obtida, nas questões discursivas do Componente de Formação Geral, por pelo menos um aluno da região Nordeste ( 92,0 ), enquanto a menor nota máxima foi obtida na região Sudeste ( 88,0 ).
Em todas as regiões a nota mínima foi zero, sem exceção.
O Gráfico 3.15 representa a distribuição das notas nas questões discursivas no Componente de Formação Geral.
Devido à grande quantidade de notas zero e à alta frequência de alunos que deixaram este tipo de questão em branco o intervalo [ 0 ; 10 ] se apresenta -se como a moda da distribuição.
Destaca -se, também, o intervalo ( 50 ; 60 ] aparecendo como uma segunda moda e com um valor próximo ao da moda principal.
A distribuição possui assimetria à esquerda, coeficiente de assimetria -0,38.
Em todas as regiões o coeficiente de assimetria também é negativo, o maior em módulo, -0,51, na região Sul, e o menor na região Sudeste, -0,23.
A avaliação das questões discursivas de Formação Geral considerou em separado o conteúdo ( peso 0,8 na nota ) e o desempenho em língua portuguesa ( peso 0,2 ).
Na sequência, os resultados verificados para cada uma das questões discursivas de Formação Geral com relação ao conteúdo serão apresentados, estabelecendo -se relações com os temas abordados em cada uma delas.
Os comentários da Banca de docentes corretores a respeito do observado na correção das respostas dos estudantes, suas impressões e conclusões serão apresentados junto à análise de cada questão.
Em seguida será feita uma análise do desempenho de Língua Portuguesa.
Os comentários da Banca de docentes corretores serão apresentados para o conjunto de questões.
Cumpre esclarecer que, tendo em vista que as questões discursivas de Formação Geral são padronizadas, ou seja, constam de todas as provas, os comentários da Banca são os mesmos para todas as carreiras acadêmicas, sendo direcionados a todos os estudantes que participaram do ENADE/2013.
A seguir, serão analisados os desempenhos da Área de Serviço Social nas duas questões discursivas de Formação Geral do ENADE/2013, comparando os resultados obtidos com comentários para cada questão. 3.3.1.1 Análise de Conteúdo Questão Discursiva 1 do Componente de Formação Geral Os dados de Serviço Social, obtidos a partir das respostas à questão 1, encontram -se na Tabela 3.9 e no Gráfico 3.16.
Nessa questão – de melhor desempenho dentre as duas de Formação Geral – os alunos de todo Brasil tiveram média, 37,7.
A maior média para a questão 1 foi obtida na região Sul ( 39,9 ), e a menor, na região Sudeste ( 34,6 ).
Quanto à variabilidade das notas, o desvio padrão de todo o Brasil foi 26,3.
O maior desvio padrão foi obtido na região Norte ( 28,0 ) e o menor desvio padrão foi obtido na região Sul ( 24,7 ).
A mediana da questão discursiva 1 foi 40,0 para todas as regiões do Brasil, exceto a região Sul ( 45,0 ).
Para o Brasil como um todo foi também 40,0.
As notas máximas ( 100,0 ) foram as mesmas para todas as regiões do Brasil, sem exceção.
Além disso, a nota mínima foi zero em todas as regiões do país, sem exceção.
O Gráfico 3.16 mostra a distribuição das notas na questão discursiva 1 do Componente de Formação Geral.
Observa -se que a maior frequência corresponde aos alunos que deixaram a questão em branco.
Considerando apenas os estudantes que tentaram resolver a questão, a moda ocorre no intervalo ( 40 ; 50 ].
A distribuição das notas possui assimetria negativa ( -0,18 ), assim como em todas as regiões do Brasil ( entre -0,26 para o Sul e -0,03 para o Norte ). 3.3.1.2 Comentários sobre a correção de Conteúdo das respostas à Questão Discursiva 1 O enunciado restou claro, elaborado com linguagem técnica, estando inserido nas Diretrizes gerais, com nível médio de profundidade, já que abordava tema vinculado à Área de formação da maior parte dos examinandos.
A divisão das perguntas também ficou clara indicando, passo a passo, o caminho de resposta a seguir, com base no texto.
Verificou -se dificuldade quanto à análise dos programas de governo existentes que contribuem para melhoria do saneamento básico.
Nesse ponto, os alunos referiram, com muita frequência, a ações relacionadas à Saúde, tais como : médico da família, ações educativas, ações ambientais, melhoria nas condições de habitação ( sem necessariamente vincular a programas de governo federal ou estaduais ou municipais ).
Os programas citados foram : Programa Saúde da Família ( PSF ), Programa de Aceleração do Crescimento ( PAC ) e Minha Casa Minha Vida e, ainda, o SUS.
Apesar de uma minoria de examinandos ter indicado o PAC do Governo Federal, a maior parte deles apresentou alguma solução para a questão do saneamento básico.
Os estudantes indicaram, com facilidade, a questão da necessidade do tratamento da água como item essencial do saneamento básico.
Muitos concluintes, através da sua experiência local, citaram situações em que as crianças caminhavam de pés no chão perto de rotas de esgoto a céu aberto.
Também criticaram a inércia governamental em não combater as invasões ilícitas, em locais inadequados, sem estrutura de água potável e também de saneamento básico.
Diante das observações preliminares pode ser extraída uma classificação das respostas obtidas que ficou estabelecida em fracas, médias e boas.
As notas consideradas fracas ( notas de 0 a 3,0 ) foram conferidas àqueles que não conseguiram relacionar o saneamento básico com a importância do tratamento de água, os projetos para resolver o problema e nem indicaram soluções para a questão.
As respostas que tiveram média avaliação ( notas de 3,5 a 7,0 ) tocaram, de algum modo, nos pontos indicados, mesmo que de forma insuficiente.
Já as respostas consideradas boas ( notas de 7,5 a 10,0 ) preencheram os requisitos do padrão de resposta, não logrando nota máxima, na maior parte dos casos, por ausência de indicação de programa governamental, mas sugerindo soluções técnicas para o problema.
As respostas mais comuns foram as consideradas medianas e que indicaram, com frequência, que o problema do saneamento básico estaria nas classes de renda mais desfavorecida e nas regiões consideradas como periferias das grandes cidades.
Também foram frequentes respostas que citaram a necessidade de tratamento de água, acompanhada da necessidade de maiores investimentos nessa área. 3.3.1.3 Análise de Conteúdo da Questão Discursiva 2 do Componente de Formação Geral A Tabela 3.10 mostra que o desempenho dos estudantes na questão 2 ( média 33,2 ) foi inferior ao obtido na questão de número 1 ( média 37,7 ).
A região Sul foi aquela onde a média, nessa questão, foi maior ( 34,5 ), e a de menor média foi a região Sudeste ( 31,0 ).
Quanto à variabilidade das notas, o desvio padrão de todo o Brasil foi 27,0, superior ao obtido na questão de número 1 ( 26,3 ).
O maior desvio nessa questão foi obtido na região Nordeste ( 29,3 ), enquanto o menor foi obtido na região Sul ( 25,3 ).
A mediana foi 35,0 em quase todas as regiões, exceto nas regiões Norte ( 37,5 ) e Sudeste ( 30,0 ).
As notas máximas ( 100,0 ) foram as mesmas em todas as regiões do Brasil, assim como as notas mínimas ( 0,0 ), sem exceção.
O Gráfico 3.17 mostra a distribuição das notas na questão discursiva 2 do Componente de Formação Geral.
Neste gráfico, desconsiderando -se a grande quantidade de alunos que deixaram a questão 2 em branco ( em torno de 24 \% ), observa -se a moda no intervalo ( 40 ; 50 ].
O coeficiente de assimetria das notas nesta questão é positivo ( 0,12 ), indicando que a distribuição é assimétrica à direita.
O mesmo ocorre para todas as regiões da Brasil ( entre 0,06 nas regiões Norte e Sul, e 0,22 na Sudeste ). 3.3.1.4 Comentários sobre a correção de Conteúdo das respostas à Questão Discursiva 2 O tema abordado no enunciado, apesar da complexidade, é extremamente atual.
Essa atualidade facilitou a elaboração de alguma resposta, permitindo que todos os que se dispuseram a responder pudessem apresentar, pelo menos, alguns dos elementos integrantes do padrão de resposta.
A correção, seguindo o padrão de resposta, avaliou a capacidade de elaborar um texto dissertativo que abordasse o tema sob três pontos de vista : o papel da tecnologia digital na atualidade ; a garantia dos direitos do cidadão e do Estado diante do avanço tecnológico ; e o problema da segurança digital para o indivíduo, para o Estado e para as organizações ( empresas e/ou instituições nacionais ou internacionais ).
Dada a amplitude de respostas possíveis, os alunos tiveram desempenho satisfatório.
Identificamos que os concluintes, basicamente, trilharam dois caminhos : a ) a análise da influência das inovações tecnológicas, notadamente através das redes sociais, numa perspectiva essencialmente privada, daí defluindo as questões da privacidade e da segurança ; b ) uma perspectiva de segurança nacional-estatal com criticas à atuação americana de invasão de privacidade dos dados governamentais, e à inercia do governo brasileiro em reagir ao ataque, considerado, pela maioria esmagadora dos concluintes, como indevido e antiético.
Poucos foram os que uniram as duas facetas da questão : o público e o privado.
Na maior parte das provas, o concluinte indicou a ausência de segurança e de privacidade decorrentes das próprias atividades individuais.
Muitos respondentes citaram que, de forma voluntária, os próprios cidadãos enfraquecem a proteção à privacidade, pois divulgam sua vida pessoal, aspectos de seu cotidiano, por meio de imagens e do fornecimento de dados pessoais.
No aspecto segurança, foram inúmeras as indicações de acessos indevidos a dados pessoais com a perpetração de fraudes de todos os tipos, notadamente as bancárias e comerciais, com a aquisição de bens à revelia dos titulares dos dados.
Também no tema segurança, muitos indicaram o aspecto positivo decorrente da vigilância generalizada, não somente através da rede de computadores, mas também pela quantidade de câmeras existentes, que contribuem para a identificação de criminosos procurados ou dos que realizam crimes pela internet.
Na perspectiva pública, a esmagadora maioria indicou repulsa pelo comportamento americano de invasão dos dados do governo brasileiro, reclamando reação do mesmo.
Alguns poucos justificaram a intervenção americana pela “ natural ” disputa de poder pelos Estados imperialistas. 3.3.1.5 Análise de Língua Portuguesa das Questões Discursivas do Componente de Formação Geral Os dados de Serviço Social, obtidos a partir das respostas às questões discursivas do Componente de Formação Geral no que tange à Língua Portuguesa, encontram -se na Tabela 3.11 e no Gráfico 3.18.
Nesse aspecto, os alunos de todo Brasil tiveram média, 40,4.
A maior média com respeito à Língua Portuguesa foi obtida na região Sul ( 43,4 ), e a menor, na região Sudeste ( 37,6 ).
Quanto à variabilidade das notas, o desvio padrão de todo o Brasil foi 23,0.
O menor desvio padrão foi obtido na região Sul ( 20,5 ) e o maior desvio padrão foi obtido na região Nordeste ( 25,6 ).
A mediana da nota de Língua Portuguesa foi 47,5 para as regiões Norte, Sul e Centro-Oeste, e mais baixa para as regiões Nordeste e Sudeste ( 45,0 ).
Para o Brasil como um todo foi também 47,5.
A nota máxima foi obtida por pelo menos um aluno tanto na região Norte quanto na região Nordeste ( 92,5 ), enquanto a menor nota máxima foi obtida na região O Gráfico 3.18 mostra a distribuição das notas de Língua Portuguesa do Componente de Formação Geral.
Observa -se que a maior frequência corresponde aos alunos que obtiveram nota no intervalo ( 50 ; 60 ], correspondendo à moda da distribuição.
Alunos que deixaram a questão em branco constituem uma moda secundária.
A distribuição das notas possui assimetria negativa ( -0,62 ), assim como em todas as regiões do Brasil ( entre -0,86 para a Sul e -0,34 para a Nordeste ). 3.3.1.6 Comentários sobre a correção das respostas de Formação Geral com respeito à Língua Portuguesa Um aspecto interessante a ser destacado é a diferença de desempenho entre a questão 1 e a questão 2, principalmente nos aspectos textuais – os textos produzidos para a questão 1 apresentaram resultados melhores do que os produzidos para a questão 2.
Uma hipótese que provavelmente explique essa discrepância é o fato de o tema “ saneamento básico ” ser mais relevante para profissionais de saúde do que o tema “ segurança e privacidade ”, constituindo tema trabalhado no currículo escolar.
Prova disso é que quase todos os textos da questão 1 citam termos específicos da área de saúde, como “ contaminação feco-oral ”, “ higiene corporal ”, “ ingestão de alimentos ou água contaminada por vetores biológicos ”.
Os resultados da avaliação correspondem aos seguintes aspectos observados em cada competência : a ) aspectos ortográficos : o desempenho dos participantes ficou concentrado, majoritariamente, nos níveis 3 ( até 3 desvios ) e 2 ( de 4 a 7 desvios ), devido ao grande índice de desvios de acentuação, contrabalançada pelo pequeno número de desvios em grafia.
Foi reduzido o número de textos com ausência de desvios ( enquadrados no nível 4 ).
Alguns textos, com pior desempenho, foram enquadrados no nível 1 ( de 8 a 12 desvios ), enquanto o nível zero foi atribuído a pouquíssimos casos, já que ele revela falta absoluta de domínio das convenções ortográficas.
Observou -se, portanto, que existe uma diferença muito grande de desempenho nos dois aspectos analisados : baixíssimo índice de desvios ortográficos e grande índice de desvios de acentuação.
Em vários casos, ocorre ausência completa de acentuação gráfica.
Os resultados revelam, portanto, que a tendência dominante entre os universitários brasileiros é a eliminação da acentuação gráfica, talvez motivada pela ausência de esclarecimento dos meios de comunicação, das autoridades e das escolas sobre as decisões do Acordo Ortográfico da Língua Portuguesa de 1990.
Os casos mais sistemáticos de ausência de acentuação são :  palavras proparoxítonas ( “ politicas ”, “ publicas ”, “ juridicas ” ) ;  palavras paroxítonas terminadas em ditongo crescente ( “ varios ”, “ Brasilia ”, “ familia ” ) ;  palavras oxítonas ( “ ninguem ”, “ esta ”, “ ate ” ).
Quanto ao domínio das convenções relativas à grafia das palavras, são pontuais os problemas encontrados ( “ infecciozas ”, “ obtos ”, “ hipedemia ”, “ saniamento ”, “ cituação ”, “ abitantes ” ).
Os únicos desvios recorrentes foram : a ausência de segmentação nos casos de verbo + pronome oblíquo “ se ”, combinada, eventualmente, à alteração da grafia do pronome ( “ percebese ”, “ percebece ”, “ preferece ” ) ; a grafia da forma verbal de pretérito imperfeito do subjuntivo como se fosse uma ocorrência de verbo + pronome oblíquo ( “ distribui -se ” por “ distribuísse ” ) Vale observar, também, que, ao contrário do que se esperava, não apareceram abreviaturas próprias do “ internetês ”, ou seja, dos hábitos de comunicação escrita adquiridos pelo uso de redes sociais e e-mails. b ) aspectos textuais : o desempenho dos participantes ficou concentrado, majoritariamente, nos níveis 3 e 2, devido à grande ocorrência de problemas de estruturação textual.
Foi muito baixo o número de textos enquadrados no nível 4 por não apresentar qualquer problema estrutural.
Alguns textos com pior desempenho foram enquadrados no nível 1, enquanto o nível zero foi atribuído a pouquíssimos casos, já que ele Observou -se que a grande maioria dos participantes não distribuiu as ideias em parágrafos, talvez devido ao pequeno número de linhas disponibilizadas para a resposta da questão ou, quem sabe, pela suposição de que não seria necessária essa divisão por não se tratar de um texto no modelo de uma redação dissertativo-argumentativa, como solicitado nos vestibulares.
Essa competência é a que se revela como a mais problemática entre os participantes, porque são muitos os problemas observados, desvios acumulados durante toda a formação do estudante e que não se resolvem com um estudo autodidata, como acontece com regras ortográficas ou morfossintáticas : sequência justaposta de ideias sem encaixamentos sintáticos ; redução drástica de estruturas subordinadas, ao lado do aumento na frequência de estruturas coordenadas e absolutas ; redução no uso de conectores para expressar relações lógicas essenciais à construção do texto, substituídas pela exigência de inferência por parte do interlocutor para suprir a sua ausência ; emprego equivocado de operadores que não estabelecem relações lógicas coerentes entre ideias do texto ; emprego inadequado do pronome relativo ( com omissão da preposição ou a utilização de pronome inadequado, como “ onde ” ) ; repetição exaustiva de termos sem a utilização de procedimentos mais sofisticados de substituição ( hiperonímias, hiponímias, nominalizações, expressões metafóricas ) ; frases fragmentadas que comprometem a estrutura lógico-gramatical ; frases formadas apenas por oração subordinada, sem oração principal.
Um importante aspecto a destacar é o baixíssimo desempenho de uma parte dos participantes em relação à estrutura formal do texto produzido, o que é extremamente preocupante ao se levar em conta que são graduandos em fase final de formação.
São frequentes os casos de desvios de estruturação frasal, com uso inadequado ou ausência de conectivos entre parágrafos e entre frases.
Em uma parte dos textos, falta um mínimo de textualidade e de domínio do registro padrão da língua.
Na verdade, observam -se relações linguísticas quase agramaticais, como as estabelecidas pela sequência de gerúndios sem o apoio de um ponto de partida para a organização das informações gramaticais e semânticas.
Quanto à utilização dos mecanismos de referenciação, dois aspectos devem ser destacados : a ocorrência, em uma boa parte dos textos, de repetições de palavras ou expressões sem a utilização de termos sinônimos ou pronomes, como seria adequado ; e a alternância entre os pronomes demonstrativos este / esse, e suas formas flexionadas, sem a observância de um critério de utilização.
Esse último aspecto tem se generalizado na modalidade oral da língua e, cada vez mais, tem se tornado uma constante na modalidade escrita de jornais e revistas : predominância da forma “ esse ” na modalidade oral e da forma “ este ” na modalidade escrita, provavelmente por um processo de hipercorreção. c ) aspectos morfossintáticos e vocabulares : o desempenho dos participantes ficou concentrado, majoritariamente, nos níveis 3 e 2, devido à grande ocorrência de problemas de regência, concordância e uso adequado dos sinais de pontuação.
O nível 4 foi atribuído a um número muito reduzido de textos.
Alguns textos com pior desempenho foram enquadrados no nível 1, enquanto o nível zero foi atribuído a pouquíssimos casos, já que ele revela a existência de um texto sem articulação e com comprometimento do sentido.
Os resultados são muito transparentes em relação aos aspectos mais problemáticos no desempenho dos participantes.
O desvio mais frequente é a falta do sinal indicativo da crase, o que se configura como um problema de regência, verbal ou nominal, conforme o caso – isso revela que o usuário não tem consciência de que, sob a forma do termo “ a ” existe a presença de uma preposição “ a ”, exigida pela regência do termo anterior.
Embora em outros exames, como o Enem, a falta de crase seja apenada em acentuação, nesta correção amostral esse desvio foi considerado no âmbito dos aspectos morfossintáticos.
Outro problema relacionado à regência verbal e à nominal, encontrado frequentemente nas questões, foi a ausência de preposição antes de pronome relativo, processo generalizado na modalidade oral da língua, em situações de registro informal.
Apesar da possibilidade de que essa alteração de regência se generalize no padrão escrito da Língua Portuguesa, como já está ocorrendo até em textos jornalísticos, o não emprego da preposição foi apenado neste processo de avaliação.
Outro desvio muito frequente diz respeito aos processos de concordância verbal e de concordância nominal.
Quanto à concordância de número, observou -se ausência de marca ( com sujeito anteposto ou posposto ) ou uso indevido ( uso inadequado da marca de plural comandado pelo núcleo plural da locução adjetiva, apesar de o substantivo que funciona como núcleo do sintagma nominal estar no singular ).
Uma ocorrência que se destacou foi a ausência de acento circunflexo na forma plural do presente do indicativo do verbo “ ter ”, que foi considerada como um desvio na concordância verbal e não na acentuação gráfica.
Quanto à concordância de gênero, vários casos foram observados, normalmente no âmbito de sintagmas nominais longos, em que o adjetivo está afastado do substantivo.
Quanto à questão da colocação pronominal, foram poucos os casos observados.
Apesar de serem aspectos relacionados à oralidade, concluiu -se que, no registro escrito formal, a maioria dos participantes já incorporou regras como a não introdução da frase por um pronome oblíquo e a próclise na presença de um termo atrator.
Não se adotou, entretanto, o padrão excessivamente formal descrito pelas gramáticas normativas em relação à posição do pronome oblíquo em locuções verbais, já que esse uso está muito distante da prática cotidiana, até em textos mais formais.
Quanto aos aspectos vocabulares, alguns tipos de inadequação foram observados : expressões da oralidade apareceram em algumas respostas, mas sem maior relevância do ponto de vista quantitativo ; seleção vocabular incompatível com o contexto, gerando situações de falta de inteligibilidade ; falta de domínio de vocabulário mais abstrato e de maior complexidade, essencial ao desenvolvimento do texto de base dissertativa.
Quanto à utilização dos sinais de pontuação, observou -se uma grande precariedade nos textos analisados.
Não foi apenada a ausência de vírgula para destacar locuções ou adjuntos adverbiais de pequena extensão deslocados de posição na frase, por ser um uso opcional.
São os seguintes os tipos de problemas encontrados : a ) vírgula : utilização de vírgula para separar o sujeito e o predicado ; ocorrência de apenas uma das vírgulas para separar uma palavra, uma expressão ou uma oração encaixada ; uso de vírgula no lugar do ponto para separar ideias que constituem períodos distintos ; ausência de vírgula para separar enumerações ; ausência de vírgula para separar oração adjetiva explicativa ou utilização inadequada para separar oração adjetiva restritiva ; b ) ponto e vírgula : utilização do ponto e vírgula no lugar de vírgula ; c ) ponto final : ausência de ponto final para separar períodos.
Várias marcas de oralidade foram identificadas, embora não em alta frequência : o uso do pronome relativo “ onde ” como relativo universal, falta de artigo definido antes de substantivo, repetição de palavras por falta de vocabulário, reduções como “ tá ”, “ pra ”, “ pro ”, “ prum ”, expressões informais como “ estão nem aí ”, “ tem que ”, “ levando com a barriga ”, “ não tá nem aí ”, “ a coisa fica séria ”, “ e aí haja dor de cabeça ”. 3.3.2 Componente de Conhecimento Específico Na parte da prova relativa às questões discursivas no Componente de Conhecimento Específico ( Tabela 3.12 ), observa -se que a média foi mais baixa do que nas questões discursivas do Componente de Formação Geral.
Enquanto no Componente de Formação Geral a média para estudantes de Serviço Social de todo o Brasil foi 36,4, na Parte de Conhecimento Específico a média foi 10,4.
A maior média deste componente foi obtida pelos estudantes da região Norte ( 13,6 ), e a menor, pelos da região Centro-Oeste ( 8,8 ).
Quanto à variabilidade das notas, o desvio padrão de todo o Brasil foi 12,8.
O maior desvio padrão foi encontrado na região Norte ( 15,7 ), e o menor, nas regiões Sul e Centro-Oeste ( 11,5 ).
A maior nota máxima, nas questões discursivas do Componente de Conhecimento Específico, foi obtida na região Nordeste ( 91,7 ), enquanto a menor nota máxima foi encontrada na região Centro-Oeste ( 73,3 ).
Além disso, a nota mínima ( 0,0 ) foi obtida por alunos de todas as regiões do Brasil sem exceção.
A mediana do Brasil como um todo foi 5,0, sendo a maior mediana encontrada nas regiões Norte e Nordeste ( 8,3 ) e a menor nas regiões Sudeste e Centro-Oeste ( 3,3 ).
O Gráfico 3.19 representa a distribuição das notas nas questões discursivas no Componente de Conhecimento Específico.
Essa distribuição tem moda no intervalo de [ 0 ; 10 ].
Nota -se que há uma tendência decrescente a partir deste intervalo até o último intervalo.
O coeficiente de assimetria é positivo para todo o Brasil ( 1,48 ) e para todas as regiões ( entre 1,23 para a região Nordeste e 1,76 para a região Centro-Oeste ), indicando uma assimetria à direita. 3.3.2.1 Análise da Questão Discursiva 3 do Componente de Conhecimento Específico Na questão 3, cujos resultados aferidos encontram -se descritos na Tabela 3.13, a média dos estudantes de todo o Brasil foi 5,1.
O desempenho dos estudantes de todo o Brasil nesta questão foi inferior ao desempenho nas questões 4 e 5.
A menor média nessa questão foi obtida pelos alunos da região Sul ( 3,8 ), enquanto a maior média foi obtida na região Norte ( 9,7 ).
Quanto à variabilidade das notas, o desvio padrão de todo o Brasil foi 11,1.
O maior desvio padrão foi obtido na região Norte ( 16,3 ), enquanto o menor foi obtido na região Sul ( 8,6 ).
A nota máxima foi alcançada por pelo menos um aluno tanto na região Nordeste quanto na região Sudeste ( 100,0 ), enquanto a menor nota máxima foi obtida na região Sul ( 80,0 ).
A mediana e a nota mínima do Brasil como um todo foi zero, o mesmo valor para todas as regiões, sem exceção, indicando que pelo menos metade dos estudantes de cada região recebeu nota zero nesta questão, seja por terem deixado a questão em branco, seja pela inadequação da resposta.
O Gráfico 3.20 mostra a distribuição das notas na questão discursiva 3, do Componente de Conhecimento Específico.
Essa distribuição é unimodal com moda nas questões com nota zero.
O coeficiente de assimetria é positivo para todo o Brasil ( 3,17 ) e para todas as regiões ( entre 2,21 na região Norte e 3,47 na região Sul ), indicando uma assimetria acentuada à direita. 3.3.2.2 Comentários sobre as respostas à Questão Discursiva 3 A Questão 3 tinha um enunciado claro e objetivo, redigido em linguagem compreensível e continha todos os elementos necessários para a compreensão da pergunta e para o desenvolvimento da resposta.
O tema era pertinente no que se refere aos conteúdos básicos previstos nas Diretrizes Curriculares - mais precisamente relacionados ao Núcleo de Fundamentos Teórico-Metodológicos da Vida Social.
O nível de profundidade de conhecimento exigido na questão era alto tendo em vista as delimitações de espaço e de tempo da prova.
No entanto, a charge utilizada como um elemento motivador para a questão, nesse caso específico, pode ter causado confusão quanto ao conteúdo a ser abordado na resposta.
De fato, foi possível verificar que houve uma incidência significativa de tentativas de respostas que procuraram estabelecer comparações - equivocadamente - entre"trabalho escravo"e"trabalho assalariado ”, o que não era o foco.
A questão solicitava a abordagem do trabalho assalariado, no modo de produção capitalista, tratando dos direitos sociais e de sua precarização, segundo o padrão de resposta.
Ainda pelo padrão, a perspectiva deveria ser a da teoria crítica social marxista : contradição fundamental do sistema capitalista ( relação capital/trabalho ) ; processo de exploração do trabalhador e extração da mais-valia ; expressões da questão social, mazelas sociais ; luta política dos trabalhadores ; direitos e precarização dos direitos.
Assim, o padrão delimitava, clara e especificamente, os eixos e categorias teóricas a serem contemplados nas respostas, permitindo relativa flexibilidade em sua construção.
Observou -se que a maioria das respostas dos estudantes não abordava os conteúdos esperados.
A interpretação da questão foi realizada de forma equivocada e a maioria das respostas era vaga, mal redigida, sem encadeamento lógico das ideias ou qualquer fundamento teórico.
Em sua maioria, as respostas se sustentavam na repetição do enunciado, sem acréscimo de qualquer elemento conceitual ou argumentativo, configurando -se, muitas vezes, em um mero rearranjo e jogo de palavras.
Outra tendência marcante foi a exposição de ideias baseadas numa quase crônica da vida pessoal do estudante perante o mercado de trabalho, sem qualquer ilação teórico-conceitual.
Uma análise global leva à percepção da dificuldade, da maioria dos estudantes, para a reflexão teórica, prevalecendo um quase-diário fenomênico, descritivo, distante do pensamento abstrato, tão necessário ao profissional da área e a essa questão discursiva, em particular.
Acrescentam -se, ainda, as dificuldades com a escrita e uso da norma culta.
A grande maioria demonstrou não ter domínio básico da língua portuguesa.
As respostas tinham problemas de concordância, pontuação, gramática, acentuação, entre outros, que comprometiam, inclusive, a clareza, objetividade das ideias e possível argumentação presente nas respostas.
Entre as respostas fracas, destacam -se aquelas que : 1 ) realizaram uma mera repetição do enunciado e/ou “ tradução ” da charge apresentada como ilustração, sem acréscimo de conteúdo ou reflexão além do já explicitado no enunciado da questão ; em alguns casos, essa repetição era acrescida de impressionismo sobre suas consequências sem, no entanto, apresentar quaisquer conceitos adicionais ; 2 ) se equivocaram na intepretação do enunciado - comparando, de forma imprecisa e inadvertidamente, o"trabalho escravo"com o"trabalho assalariado"; alguns, inclusive, pareciam fazer a"defesa"da"superioridade"da produção escravocrata como um sistema mais justo e digno do que o do modo de produção capitalista ; 3 ) houve relatos de casos e/ou denúncias de situações de trabalho escravo, com base no senso comum ; 4 ) com maior frequência verificou -se grande dificuldades de redação, ausência de pensamento organizado, fuga de ideias e exploração sentimental do tema.
De fato, aspectos dos processos sociais históricos foram, muitas vezes, deturpados e estavam presentes em expressões sobre “ a revolução industrial no Brasil ”, ou “ os escravos no feudalismo... ”, ou, ainda, na ausência de qualquer datação histórica, ao se comparar a realidade atual com"algumas décadas atrás"para caracterizar, cronologicamente, o escravismo no país.
Sobretudo, vale ressaltar o renitente desconhecimento da formação histórica e social do Brasil e, ainda maior, do mundo.
A ausência de uma cultura mais abrangente, e retrospectiva, provocou respostas errôneas sobre o processo escravocrata no Brasil Colônia e Império e também sobre o modo de produção capitalista.
Foram observadas, também, respostas que reportaram toda a história das lutas sociais dos trabalhadores como resumida e/ou sintetizada nos direitos conquistados na Constituição Federal de 1988, ou que afirmavam que todos os processos de dominação capitalista teriam sido engendrados pelo neoliberalismo.
Foi possível identificar o hiperdimensionamento de fatores como a fragilidade da qualificação dos trabalhadores e o perfil da “ mão de obra ” hoje.
A educação, a qualificação e/ou a capacitação dos trabalhadores apareceram como a questão-chave das péssimas condições de trabalho, da precarização e do desemprego.
De certo modo, variadas vezes, os meios de comunicação foram citados como referências analíticas, sem qualquer ponderação crítica dessas fontes.
Nesses exemplares de respostas foi possível evidenciar a tendência à reprodução do senso comum na elaboração das narrativas, com ênfase em discursos utilitários e pouco analíticos.
Verificou -se, também, que as respostas, de modo geral, continham acentuada ênfase em uma concepção de história fatalista, nada próxima da teoria marxista solicitada como referência no enunciado da questão.
Esse fatalismo pode ser verificado em expressões como : “ capitalismo nunca vai acabar ” ; “ a exploração sempre existiu ” ; “ continuam sendo escravos ”.
Esse fatalismo levou a conclusões impressionistas sobre possíveis “ salvadores ” : Deus, assistente social, Estado ou fiscalização.
A propriedade privada, a divisão do trabalho, o valor e as classes sociais, de modo geral, ficaram longe das narrativas sobre a exploração dos trabalhadores.
Algumas respostas seguiam uma racionalidade de culpabilização dos trabalhadores, quando evidenciavam que “ tem exploração porque eles não sabem os direitos ” ou o “ êxodo rural ” é o causador dos problemas de sobrevivência dos trabalhadores.
Desse modo, os estudantes demonstravam que não compreendiam a sociedade capitalista enquanto totalidade social, o que é uma condição básica da compreensão da dialética marxista.
Isso foi evidenciado também quando as respostas enfatizavam dimensões parciais da reflexão de Marx sobre o trabalho e sobre as relações sociais em geral, que só podem ser compreendidas no quadro da especificidade capitalista, o que nem sempre foi abordado pelos estudantes, prevalecendo ilações parciais como : “ o trabalho dignifica o homem ” ou “ a dominação do homem pelo homem ”.
Marcas de discurso vazio por não estarem situadas nas relações sociais e de produção, configurando -se como visões a-históricas e moralistas dos processos sociais.
Como já mencionado, na verdade, a narrativa geral é “ auto-referenciada ” e, assim, as respostas se restringem a discorrer sobre a necessidade de melhoria de programas sociais, o desenvolvimento de mais políticas sociais e mais campo de trabalho para os assistentes sociais.
Uma visão corporativa que pinça, da realidade e do conhecimento acumulado, aqueles elementos que justificam a Área, o que não se coaduna com o pensamento profissional marxista no Serviço Social.
As determinações estruturais do metabolismo do capital e a ação das classes sociais ficam apagadas do universo explicativo da realidade, o que contrasta negativamente com significativa parcela da produção bibliográfica da área profissional e com as diretrizes curriculares da formação profissional do assistente social no Brasil.
Parte significativa dessas respostas obtiveram a nota 0 ( zero ).
Outras, consideradas como"respostas fracas"continham algum dos elementos esperados no padrão de resposta, apresentando -se no limiar desse quadro mais geral impressionista e a-histórico.
As respostas consideradas medianas, via de regra, atenderam parcialmente aos tópicos estabelecidos no padrão de resposta, fixando -se em parte dos conteúdos dos eixos delimitados, reportando -se, ainda que também de forma parcial, aos conteúdos teórico-conceituais na abordagem da questão.
Contudo, tais respostas não se constituíam como fruto de um desempenho efetivamente qualificado.
Isso porque, de uma maneira abrangente, com poucas exceções, as respostas medianas avançaram na melhoria do texto escrito, na sistematização do conhecimento e no tratamento mais adequado do tema, embora permanecesse o vazio da crítica e o reduzido conhecimento teórico da questão.
As boas respostas continham conteúdos pertinentes aos eixos temáticos previstos no padrão de resposta e reportavam -se aos conceitos e categorias teóricas relacionadas à questão do trabalho assalariado e da precarização dos direitos.
Há que se lamentar, todavia, a escassez de respostas com essa qualificação, uma vez que não representa, relativamente ao conjunto das respostas, um número significativo.
O elevado número de notas zero é um dado preocupante, tendo em vista que a questão aborda conceitos estruturantes do curso de graduação em Serviço Social.
Não há, portanto, muito a comemorar a respeito do desempenho dos estudantes, apesar de ser possível exemplificar respostas boas.
Cabe ainda mencionar que foram identificadas respostas que reproduziam parte de enunciados e respostas de questões objetivas - às vezes, associando mais de uma daquelas questões.
Esse tipo de recurso exigiu a atenção redobrada dos corretores e conhecimento de todo conteúdo da prova, como já é de praxe. 3.3.2.3 Análise da Questão Discursiva 4 do Componente de Conhecimento Específico A Tabela 3.14 contém as informações relativas à questão 4 do conjunto de questões do Componente de Conhecimento Específico.
O desempenho dos estudantes de todo o Brasil nesta questão foi superior ao desempenho nas questões 3 e 5.
A média geral do Brasil foi 18,1, sendo a menor média registrada na região Centro-Oeste ( 16,1 ) e a maior na região Nordeste ( 20,3 ).
Quanto à variabilidade das notas, o desvio padrão de todo o Brasil foi 24,2.
O maior desvio padrão foi obtido na região Nordeste ( 26,7 ), enquanto o menor foi obtido na região Centro-Oeste ( 22,4 ).
A nota máxima ( 100,0 ) foi atingida em todas as regiões, sem exceção.
A mediana e a nota mínima foram zero para quase todas as regiões do Brasil, exceto a mediana nas regiões Norte e Sul ( 5,0 ).
Ou seja, pelo menos metade dos alunos das regiões Nordeste, Sudeste e Centro-Oeste tirou nota zero nesta questão.
O Gráfico 3.21 representa a distribuição de notas na questão discursiva 4, no Componente de Conhecimento Específico.
Esta distribuição é unimodal com moda nas questões em branco.
O coeficiente de assimetria é positivo para todo o Brasil ( 1,16 ) e para todas as regiões ( entre 1,02 para a região Nordeste e 1,30 para a região Norte ), indicando uma assimetria à direita. 3.3.2.4 Comentários sobre as respostas à Questão Discursiva 4 O enunciado da Questão 4 era claro, redigido em linguagem acessível aos estudantes de Serviço Social e com temática adequada às Diretrizes Curriculares da ABEPSS, por abordar os aportes legais e os processos de planejamento que fazem parte de conteúdos ementários presentes no processo de formação profissional.
A questão era objetiva e sua dificuldade envolvia três momentos cruciais : a apreensão do fenômeno sobre o qual se deveria propor a intervenção a partir da leitura da tabela ; a identificação correta do que considerar em cada item da proposta ( justificativa, objetivo e metodologia ) e, por último, o correto conhecimento do que caracteriza a ação profissional.
Somam -se a isso as referências às legislações, seja aquela diretamente relacionada ao fenômeno levantado na questão, seja aquela relativa ao projeto ético-político profissional.
Cabe ressaltar que esse formato da questão exige a elaboração de um texto com níveis de profundidade e articulação difíceis de serem atendidos – em sua extensão e conteúdo – nos limites de uma prova com as características do ENADE.
A Questão 4, portanto, foi considerada bastante abrangente, com um grau elevado de exigência para os estudantes, considerando -se a extensão e escopo da prova.
A necessidade de utilização de dados apresentados em uma tabela exige : a capacidade de interpretação de situações que fazem parte do cotidiano de trabalho do assistente social e conhecimento de diferentes formas de registro e de análise, o que é pertinente e desejável.
Destaca -se que o ponto de partida para a apreensão do fenômeno seriam os elementos presentes na tabela, que apresentava distintas variáveis, com desagregações extensivas.
A partir dessa análise inicial é que seria possível apreender o fenômeno sobre o qual o aluno deveria apontar a intervenção profissional.
Tal projeto de intervenção deveria ter como base a legislação de proteção à criança e ao adolescente e o Código de Ética Profissional do Assistente Social, enfocando três de seus elementos constitutivos, que tinham pesos diferenciados em termos de pontuação.
O padrão de resposta requeria a identificação de alguns elementos constitutivos de um projeto de intervenção, o que constituiu subitens para a avaliação da questão, bem como os conteúdos centrais que deveriam constar em cada tópico do projeto.
Nesta direção, o padrão de resposta, coerente com o enunciado, determinava o que cada item deveria conter em termos de conteúdo.
Entretanto, o conjunto de elementos destacados nos três itens era demasiadamente extenso, considerando o curto espaço de linhas e de tempo que o estudante teve para compreender o enunciado da questão, analisar os dados da tabela, e elaborar a resposta de cada item.
Apesar de extenso, o conteúdo exigido é compatível com a formação profissional do Assistente Social e permite avaliar o grau de compreensão dos estudantes acerca das razões que justificam uma ação profissional sobre a situação de violência apresentada, com base na utilização de instrumental técnico.
O padrão permitiu avaliar se o estudante possuía um conhecimento básico sobre o que é necessário para elaboração de um projeto de intervenção.
Todavia, foi observado um grande distanciamento entre o esperado pelo padrão resposta e as respostas dadas pelos estudantes que, em grande medida, veiculavam conteúdos panfletários e moralistas, sem se ater à definição dos elementos do projeto.
Portanto, foram frequentes as respostas vagas e superficiais elaboradas em texto corrido, sem subdivisão, que não atendiam ao solicitado, e que continham variadas interpretações dos dados apresentados na tabela e também sobre o que era exigido na questão.
Foi possível observar que a grande maioria das respostas não abordou os conteúdos relacionados aos itens exigidos.
Como a questão solicitava um projeto de intervenção, no geral as respostas eram apresentadas sem a correspondência com os itens do projeto, quais sejam, justificativa, objetivo e metodologia.
Além disso, mesmo as respostas que possuíam a definição dos itens ‘ a ’, ‘ b ’ e ‘ c ’, correspondentes aos respectivos elementos do projeto, apresentavam um conteúdo que não se relacionava ao exigido e ao esperado. quais sejam, justificativa, objetivo e metodologia.
Além disso, mesmo as respostas que possuíam a definição dos itens ‘ a ’, ‘ b ’ e ‘ c ’, correspondentes aos respectivos elementos do projeto, apresentavam um conteúdo que não se relacionava ao exigido e ao esperado.
Poucos alunos atenderam plenamente ao esperado no padrão de resposta.
As respostas dadas pelos estudantes nesta questão podem ser classificadas de três formas : aqueles que, seguindo as orientações do enunciado, elaboraram proposições a partir dos itens de um projeto de intervenção ; aqueles que desenvolveram a resposta em um texto corrido, para os quais a identificação dos itens não era explícita no corpo da dissertação ; e aqueles que desenvolveram textos desarticulados, respostas vagas e sem conexões entre os elementos de um projeto, contendo abordagens sem qualquer embasamento teórico, criticidade e totalmente descontextualizadas.
A maior parte dos estudantes teve dificuldade em identificar com precisão os elementos do projeto.
Foram frequentes as confusões feitas entre projeto de intervenção e projeto de pesquisa, que se expressaram, sobretudo, no item metodologia, com proposições de técnicas e métodos de coletas de dados no lugar de ações interventivas propriamente ditas.
De um modo geral, os alunos também confundiram justificativa com objetivo e metodologia, mais uma vez, revelando desconhecer o significado dos elementos que compõem um projeto de intervenção profissional.
Foram poucas as respostas que articulavam os itens ( objetivos, justificativa e metodologia ) entre si, com coerência e unidade.
Isto é, um número reduzido de estudantes apresentou respostas estruturadas com lógica e conteúdos pertinentes, de forma fundamentada e relacionada aos termos do debate teórico-metodológico e ético-político hegemônico na profissão na atualidade.
As respostas relativas ao item ‘ justificativa ’ continham concepções ideologizadas acerca da temática ou pautadas no senso comum e pouco se reportavam à legislação referente aos direitos das crianças e adolescentes e ao exercício profissional, tais como o Código de Ética do Assistente Social e a Lei de Regulamentação da Profissão.
No item referente ao ‘ objetivo ’ as indicações pouco se aproximaram do padrão de resposta, prevalecendo formulações para além da capacidade de atuação profissional ou compostas por conteúdo moral ou educativo que em nada guarda relação com a ação técnica, política e teórica do assistente social.
No que diz respeito ao item ‘ metodologia ’ as proposições técnicas e de com quem deveriam ser realizadas as ações foram aquelas que predominaram e elevaram as notas.
No conjunto das respostas, foi possível observar o predomínio de uma abordagem que privilegiou o item de metodologia através de indicações de possibilidades de realização de atividades com as “ vítimas de violência familiar ”.
Nesse item, os estudantes apresentaram maior facilidade em responder, recorrendo a propostas de modo instrumental.
Cabe ressaltar, entretanto, que foram formuladas sugestões pontuais, soltas e sem articulação com os demais itens em sua grande maioria.
Observou -se, também, que um conjunto significativo de respostas fazia referência a situações hipotéticas de violência doméstica, em municípios variados, sem diretamente vincularem o projeto de intervenção proposto aos indicadores apresentados no enunciado da questão.
Apesar disso, quando as respostas apresentavam desdobramentos em termos de propostas de intervenção, foram consideradas.
Muitas vezes, as respostas dos estudantes se limitavam a opinar sobre a questão da violência intrafamiliar contra as crianças, sugerindo medidas preventivas, de forte diretriz moralizadora, ou medidas ex post facto, voltadas para os danos causados às vítimas.
Nesse viés, a intervenção do Assistente Social se materializa na habilidade de conduzir a criança a revelar a violência sofrida, para em seguida encaminhar o caso à autoridade, pedir o afastamento e a punição do culpado.
A perspectiva da punição parece ser o objetivo primordial da ação profissional, ignorando o fato de a violência doméstica ocorrer no contexto de relações significativas para a criança.
O uso de nomenclatura já superada pelo Serviço Social também apareceu com frequência, prevalecendo o Assistente Social como o profissional que detém um saber capaz de “ diagnosticar ” a violência doméstica e promover a punição do culpado.
O discurso prevalente demonstra que o aluno confundiu as atribuições do Assistente Social com as funções de julgar, punir e corrigir.
As respostas, frequentemente, expressavam uma compreensão ingênua da família, que aparece como culpada ou vítima da sua própria transformação.
O uso de álcool e outras drogas, o trabalho da mulher e sua consequente ausência de casa, bem como a recomposição das famílias, são apontados como causa em si, da violência intrafamiliar, ignorando o dever do Estado e o papel das políticas públicas.
Soma -se a isso que, comumente, a questão da violência sexual foi tratada com um misto de preconceito, moralismo e horror, sem a compreensão que o tema exige, dada a sua complexidade.
Pareceu -nos que a apreensão ingênua do fenômeno conduziu a interpretações equivocadas, justificando a abordagem conservadora.
Finalmente, cabe considerar que um conjunto significativo de questões relacionou a violência doméstica à “ desestruturação familiar ” ou à “ falta do amor a Deus ”.
No caso da tese da desestruturação familiar, algumas respostas apontavam ações que colocam a família como central, propondo trabalhos em grupos, atendimento terapêutico individual e etc.
Acreditamos que a dificuldade de leitura e interpretação da tabela apresentada no enunciado, que apresentava um leque variado e complexo de variáveis relacionadas às situações de violência contra crianças e adolescentes, aliada à exiguidade do tempo causaram dificuldades adicionais aos estudantes.
A grande maioria das notas foi baixa, pois não continham um projeto de intervenção profissional.
Muitas das respostas meramente apresentavam as opiniões dos alunos quanto ao tema e não a elaboração de um projeto.
Assim, a maioria das respostas não apresentou justificativa e objetivo e, quando havia algum aspecto a ser pontuado, era relativo a atividades e inerentes à metodologia.
As respostas fracas se caracterizaram por serem redações com forte conteúdo moralista e conservador que, muitas vezes, interpretam os fenômenos da violência, da violência doméstica, de maus-tratos, por exemplo, de forma preconceituosa, rasa e simplificada, em nada correspondendo ao esperado como resposta.
Dentre as respostas fracas muitas evidenciavam confusão entre a atuação do Serviço Social e a transmissão de valores morais e religiosos, contrariando os fundamentos da profissão e os princípios éticos que a regem.
Diversas expressavam uma visão conservadora da sociedade, com a proposição de medidas contrárias aos princípios básicos de respeito à pessoa e aos direitos humanos, em desagravo ao Código de Ética profissional.
Nessas redações, os estudantes se mostram incapazes de superar o senso comum, expressando uma visão moralizadora, sentimentalista e corretiva.
Predominaram nas respostas fortes conteúdos disciplinadores das relações familiares, alusivas ao resgate de valores morais, sem inserir a problemática da violência intrafamiliar no contexto social e das políticas públicas.
Também foram recorrentes as respostas com forte conteúdo religioso, baseadas em discursos psicologizantes, de responsabilização e culpabilização do indivíduo e da família ( desestruturação familiar ).
Nas respostas com notas medianas, apesar do objetivo geral se manter muito distante do que de fato compete ao assistente social, a justificativa trazia mais elementos de reflexão e a metodologia trazia ao menos dois itens.
Na metodologia foram recorrentes : o modo como o trabalho seria realizado, com clara predominância da escolha dos instrumentos e das atividades a serem realizadas, e o tipo de abordagem ( ações junto às famílias e às escolas predominam ).
Desse modo, as respostas medianas, em grande medida, tiveram a pontuação alavancada pelo item"metodologia".
Entretanto, apesar da pontuação 5,0 estar relacionada a um único item do projeto - a metodologia -, muitas vezes, as proposições feitas pelos estudantes não guardavam qualquer correspondência com o que tinha sido exposto como justificativa e objetivo.
Com certa frequência, observou -se dissociação entre a justificativa, os objetivos e a metodologia, e que os estudantes não inscreveram suas respostas no campo dos direitos sociais.
Um aspecto interessante nesse bloco de respostas foi o relevo que os alunos deram à política de Educação e à escola, considerada a melhor instituição para prevenir, identificar e enfrentar a violência contra crianças e adolescentes.
Sendo assim, os estudantes endossaram a necessidade do profissional de Serviço Social atuar na área da Educação, principalmente intervindo sobre a violência contra a criança e o adolescente.
Soma -se a isso, a visão de que a Educação é supostamente a grande solução para os diversos processos intervenientes na problemática da violência contra crianças e adolescentes.
Nessa linha de raciocínio, a demanda por creches e escolas em horário integral, o que possibilitaria que as crianças ficassem menos expostas a situações de violência.
A permanência do jovem no espaço escolar contribuiria, também, para a prevenção ao uso e abuso de drogas, para a ampliação da escolarização, mediante o reforço escolar, e mesmo para a qualificação para o mercado de trabalho, dentre outras perspectivas.
Em resumo, além da identificação da violência, a escola foi considerada o locus privilegiado para a promoção de uma adequada assistência à infância e à adolescência.
As respostas que se sustentaram no Estatuto da Criança e do Adolescente ( ECA ) para a garantia de direitos e para a apresentação de ações de enfrentamento mediadas por essa lei, quase sempre eram boas.
Tais respostas demonstravam o entendimento de que as crianças e adolescentes são sujeitos de direitos e preconizam a sua proteção integral, a partir de ações de políticas públicas intersetoriais, interdisciplinares e de acolhida das famílias nas situações de violência e na promoção de uma cultura protetiva por parte da família e da sociedade.
A maior parte das boas respostas teve como foco a promoção dos direitos das crianças e adolescentes, com destaque para a dimensão socioeducativa dos projetos e para ações de prevenção mais do que aquelas de enfrentamento quando a violência está instalada no ambiente doméstico.
Os estudantes se apoiaram nos direitos da criança e do adolescente e apontaram a necessidade de garanti -los, para justificar o projeto de intervenção.
Construíram um objetivo coerente com a justificativa, na direção da promoção destes direitos.
Na metodologia, os estudantes foram capazes de indicar os tipos de abordagem, bem como os instrumentos a serem utilizados no exercício profissional.
Ao final da correção foi possível identificar o predomínio da falta de criticidade e criatividade profissional frente a realidade social concreta e suas contradições.
Vale destacar um número significativo de respostas que indicaram ações que não cabem ao assistente social realizar, particularmente, aquelas que são de competência do Estado, da dissociação entre a justificativa, os objetivos e a metodologia, e que os estudantes não inscreveram suas respostas no campo dos direitos sociais.
Um aspecto interessante nesse bloco de respostas foi o relevo que os alunos deram à política de Educação e à escola, considerada a melhor instituição para prevenir, identificar e enfrentar a violência contra crianças e adolescentes.
Sendo assim, os estudantes endossaram a necessidade do profissional de Serviço Social atuar na área da Educação, principalmente intervindo sobre a violência contra a criança e o adolescente.
Soma -se a isso, a visão de que a Educação é supostamente a grande solução para os diversos processos intervenientes na problemática da violência contra crianças e adolescentes.
Nessa linha de raciocínio, a demanda por creches e escolas em horário integral, o que possibilitaria que as crianças ficassem menos expostas a situações de violência.
A permanência do jovem no espaço escolar contribuiria, também, para a prevenção ao uso e abuso de drogas, para a ampliação da escolarização, mediante o reforço escolar, e mesmo para a qualificação para o mercado de trabalho, dentre outras perspectivas.
Em resumo, além da identificação da violência, a escola foi considerada o locus privilegiado para a promoção de uma adequada assistência à infância e à adolescência.
As respostas que se sustentaram no Estatuto da Criança e do Adolescente ( ECA ) para a garantia de direitos e para a apresentação de ações de enfrentamento mediadas por essa lei, quase sempre eram boas.
Tais respostas demonstravam o entendimento de que as crianças e adolescentes são sujeitos de direitos e preconizam a sua proteção integral, a partir de ações de políticas públicas intersetoriais, interdisciplinares e de acolhida das famílias nas situações de violência e na promoção de uma cultura protetiva por parte da família e da sociedade.
A maior parte das boas respostas teve como foco a promoção dos direitos das crianças e adolescentes, com destaque para a dimensão socioeducativa dos projetos e para ações de prevenção mais do que aquelas de enfrentamento quando a violência está instalada no ambiente doméstico.
Os estudantes se apoiaram nos direitos da criança e do adolescente e apontaram a necessidade de garanti -los, para justificar o projeto de intervenção.
Construíram um objetivo coerente com a justificativa, na direção da promoção destes direitos.
Na metodologia, os estudantes foram capazes de indicar os tipos de abordagem, bem como os instrumentos a serem utilizados no exercício profissional.
Ao final da correção foi possível identificar o predomínio da falta de criticidade e criatividade profissional frente a realidade social concreta e suas contradições.
Vale destacar um número significativo de respostas que indicaram ações que não cabem ao assistente social realizar, particularmente, aquelas que são de competência do Estado, da justiça, e mesmo de outras áreas profissionais como a psicologia.
Todas essas observações retratam o baixo desempenho dos alunos nesse exame, evidenciando que a formação profissional em Serviço Social necessita de uma profunda revisão.
Preocupa, ainda, a baixíssima qualidade do texto apresentado pelos alunos.
Os erros rudimentares de ortografia e gramática e a péssima caligrafia exigiram um enorme esforço de leitura e compreensão por parte dos corretores. 3.3.2.5 Análise da Questão Discursiva 5 do Componente de Conhecimento Específico A Tabela 3.15 contém as informações relativas à questão 5 do conjunto do Componente de Conhecimento Específico.
O desempenho dos estudantes nessa questão foi superior ao da questão 3 e inferior ao da questão 4.
A nota média dos estudantes de todo o Brasil foi 8,1.
A maior média foi registrada na região Norte ( 12,4 ), enquanto a menor média foi registrada na região Centro-Oeste ( 6,4 ).
Quanto à variabilidade das notas, o desvio padrão dos alunos do Brasil, como um todo, foi 15,0.
Enquanto o maior desvio foi encontrado na região Norte ( 19,6 ), o menor foi encontrado na região Sul ( 13,0 ).
A nota máxima ( 100,0 ) foi alcançada em quase todas as regiões do Brasil, exceto nas regiões Nordeste e Sudeste ( 95,0 ).
A mediana e a nota mínima foram zero para todas as regiões do Brasil, sem exceção.
Ou seja, pelo menos metade dos alunos tirou nota zero nesta questão.
O Gráfico 3.22 representa a distribuição das notas da questão discursiva 5 do Componente de Conhecimento Específico.
Destaca -se, como na questão discursiva 4, o grande número de estudantes que deixaram esta questão em branco, correspondendo à moda da distribuição.
O coeficiente de assimetria para todos os alunos é 2,30 e para todas as regiões se mantém positivo, variando entre 1,66 na região Norte e 2,93 na região Centro-Oeste. 3.3.2.6 Comentários sobre as respostas à Questão Discursiva 5 O enunciado da Questão 5 apresentava uma tabela com a evolução do número de homicídios por cor/raça das vítimas ao longo dos anos de 2002 a 2010.
A questão abordava um conceito chave à compreensão do objeto de trabalho de todo e qualquer assistente social.
Como na Questão 4, a partir dos dados apresentados, os estudantes deveriam elaborar uma proposta de ação profissional.
Para a proposta a ser apresentada nessa questão o estudante deveria levar em conta, como consta do enunciado, que o tema do preconceito e da discriminação racial guarda relação com os princípios contidos no Código de Ética profissional.
O aluno deveria abordar a relação entre um dos conceitos centrais na formação profissional do assistente social, a questão social e a pobreza, de modo a contribuir para a eliminação do preconceito e da discriminação, conforme está previsto no Código de Ética Profissional.
Para isso, a questão continha elementos suficientes para orientar os alunos, em consonância com o Código de Ética Profissional.
Cabe destacar que os dados da tabela reforçam que a questão da violência tem um claro componente racial e que incide sobre a parcela da população mais pobre.
Assim, a relação entre pobreza e questão social deveria constituir o ponto de partida da resposta, requerendo que fosse explicitado o conceito de questão social, associando -o à pobreza como uma de suas expressões.
O padrão de resposta indicava a necessidade de três movimentos na correção da resposta.
Primeiramente, se deveria observar se o estudante demonstrava alguma compreensão acerca do significado do conceito de questão social e de sua vinculação com a dinâmica própria da sociedade capitalista, além de tratar a pobreza como uma de suas expressões.
Em seguida, verificava -se se a resposta continha uma abordagem dos princípios do Código de Ética e/ou do Projeto Ético-Político do Serviço Social, evidenciando uma postura crítica em relação a todas as formas de preconceito e discriminação.
A seguir, o padrão de resposta exigia a proposição de ações que assegurem esses princípios como parte das atividades que caracterizam o exercício profissional, podendo ser apontadas diferentes atividades.
O padrão de resposta da Questão 5 contemplava vários itens agrupados em duas partes, referentes a cada um dos tópicos solicitados na questão.
Cabe destacar que os conteúdos exigidos pelo padrão de resposta eram coerentes do ponto de vista teórico e não havia sobreposição, o que favoreceu a diferenciação de cada item.
O padrão de resposta cobriu diferentes tipos de respostas a partir de um eixo central.
Ao responder o item ‘ a ’, a maioria dos candidatos não articulou de forma consistente a relação entre questão social e pobreza.
Quando esta relação foi feita, apenas havia referência à pobreza como expressão da questão social sem se deter em justificar como suas expressões tem origem e se manifestam no capitalismo.
Parte significativa, ainda ao responder a primeira parte da pergunta ( relação entre a questão social e pobreza ), apenas mencionou que a pobreza é uma expressão da questão social, sem antes ter conceituado esta.
As respostas pareciam expressar um jargão já incorporado, sem que os estudantes tivessem, de fato, consciência dos meandros desta afirmação.
Outro conjunto de estudantes, respondeu parcialmente à primeira parte da pergunta e a seguir fez menção à discriminação racial como resultado do processo histórico brasileiro, referindo -se principalmente aos séculos em que a escravidão vigorou no país, justificando, a partir daí, o elevado número de homicídios e a vitimização da população negra.
Houve também um grupo de respostas que repetiu as informações contidas na tabela apresentada e as considerações feitas sobre os princípios do código de ética do Assistente Social, que condena o preconceito e a discriminação, não acrescentando conteúdo à resposta, não demonstrando, portanto, conhecimento sobre os temas solicitados. conteúdo à resposta, não demonstrando, portanto, conhecimento sobre os temas solicitados.
Ressalte -se que um pequeno número de estudantes conseguiu responder satisfatoriamente a esta primeira parte da questão, conceituando a questão social como resultante das relações capitalistas e seus impactos na produção das desigualdades, onde a pobreza e a discriminação social se inserem.
Ao responder a letra ‘ b ’, muitos estudantes apontaram propostas de ação profissional com vistas ao enfrentamento do preconceito e da discriminação racial, sendo o mais frequente a realização de palestras ( 1º lugar ) e de grupos ( 2º ação mais apontada ).
Também foi mencionada com menor frequência a realização de oficinas, debates e passeatas.
Em muitos casos, tais proposições se restringiam ao senso comum e ao voluntarismo, em face das insuficiências das respostas dadas pelos estudantes ao item ‘ a ’ ( questão social X pobreza X discriminação racial e vitimização da população negra ).
Apesar disso, foram atribuídos pontos a esta segunda parte, a despeito de sua coerência ou vinculação com a resposta dada na letra ‘ a ’.
Poucas respostas fizeram menção ou discutiram de forma consistente os dispositivos do código de ética profissional e outras legislações da profissão que tocam na questão levantada pela pergunta.
Houve também um conjunto bastante expressivo de respostas de conteúdo religioso e proposição de ações de voluntariado ressaltando a necessidade de mais “ humanidade ” ( termo muito utilizado ) frente aos problemas vivenciados pela população negra em nossa sociedade.
De forma geral, as respostas foram muito fracas.
Observou -se uma grande utilização dos termos “ problemas sociais ” e “ vulnerabilidade social ”, tomados como sinônimos de questão social e suas expressões.
Percebemos um grave desconhecimento acerca da definição de questão social, contida na bibliografia profissional, muitas vezes, descrita como instrumento ou como objetivo do trabalho do assistente social.
O tratamento dado à relação entre questão social e o trabalho profissional pautou -se, na maioria das vezes, no reconhecimento daquela como uma “ entidade ”, um “ sujeito ” capaz de interferir ou determinar o modo de atuação da profissão.
Chama a atenção, ainda, o entendimento de que o preconceito e a discriminação devem ser combatidos por meio de uma mudança de postura pessoal de base religiosa, como uma postura que deve estar presente no trabalho profissional.
Causou surpresa a percepção demonstrada por muitos alunos de que a política de cotas para negros e afrodescendentes é expressão do preconceito e da desigualdade social, devendo, para muitos, ser abolida.
Um número expressivo de respostas deixou clara a prevalência de uma visão ancorada no senso comum, revestida de conteúdo moralista e prescritivo.
O completo distanciamento dos respondentes dos aspectos exigidos : relação entre questão social e serviço social e uma proposta de ação profissional, leva a uma indagação sobre a qualidade da formação.
Muitas respostas indicaram dificuldade na interpretação do enunciado ; o estabelecimento da relação com a questão social, pelo desconhecimento de sua base conceitual, o que propiciou uma construção teoricamente frágil.
Um número representativo de respostas também apresenta grande distanciamento em relação ao tema central do preconceito racial.
Nelas os estudantes discorriam sobre outras temáticas como envelhecimento ou deficiência.
A ausência de propostas para eliminação do preconceito também foi observada, apontando para a dificuldade na construção de ações interventivas.
A expressão escrita é outro fator preocupante não só na grafia, cujos erros crassos dificultavam o entendimento de palavras, por exemplo, como “ bullying ” escrito como “ bulem ”.
A construção textual, quase sempre, era confusa.
O quantitativo elevado de documentos com nota zero é em parte sustentado pelos seguintes fatores : 1 ) não tratamento da questão de forma objetiva, ou seja, não responderam ao solicitado ou não entenderam o que fora solicitado ; 2 ) graves problemas de exposição das ideias, o que inviabilizou a construção de textos com o mínimo de sentido ; 3 ) a opção por uma linguagem coloquial, pautada no senso comum e/ou numa abordagem de cunho moralista ou religioso.
Por último, sinalizamos a desarticulação, quase sempre presente entre as duas partes da resposta.
A primeira com uma clara personificação da questão social, que deixou de ser tratada como um conceito central na compreensão da ação e história profissional para se tornar um ente definidor de distintos processos e problemas sociais.
Em relação à segunda parte, houve prevalência de ações completamente distantes do universo profissional, mesmo a indicação de palestras, políticas e outras atividades surgiam sem qualquer referência ao Serviço Social.
As respostas com boa avaliação eram as mais objetivas e, portanto, aproveitavam ao máximo o pouco espaço disponível.
Contudo, não se deve desconsiderar que a enorme quantidade de questões respondidas fora do espaço correto sinaliza algum grau de dificuldade dos alunos em organizar as ideias prévia e mentalmente, fazendo uso recorrente dos espaços destinados às questões como rascunhos de algumas respostas ou sinalizando a falta de condições de responder em função do excesso, real, de questões dissertativas, com tabelas após uma exaustiva prova objetiva. 3.3.3 Considerações Finais Os sentidos da avaliação do ENADE estão claramente definidos em seu objetivo “ avaliar o desempenho dos estudantes com relação aos conteúdos programáticos previstos nas diretrizes curriculares dos cursos de graduação, o desenvolvimento de competências e habilidades necessárias ao aprofundamento da Formação Geral e profissional, e o nível de atualização dos estudantes com relação à realidade brasileira e mundial ” ( BRASIL, 2004)12.
A busca da concretização desses sentidos se materializou na realização dos exames nacionais de avaliação de desempenho, sendo o último o de 2013, agora objeto dos presentes comentários.
A realidade observada durante as correções sequer tangencia levemente a expectativa de identificar desempenhos satisfatórios e um perfil de profissionais com pensamento crítico, autônomo e propositivo.
Ao contrário, o que se observa é um contingente de universitários perdidos entre o imediatismo da realidade vivida e a essência verdadeira dessa realidade.
E assim, transferem para as respostas das questões discursivas impressões pessoais, “ achismos ”, respostas piegas, acríticas, ou redentoras, presumidamente retiradas da realidade imediatamente percebida pelos autores das respostas.
E, de fato, a prevalência de notas baixas espelha as dificuldades dos estudantes, que participaram do ENADE, para atender ao que era esperado pelo padrão de respostas.
Assim, é possível inferir dois problemas estruturais básicos.
De um lado, fragilidade formativa em termos teórico-conceituais, provocando sérias limitações na capacidade de generalização e de abstração sobre o concreto.
De outro lado, essa dificuldade de afastamento da mera descrição de experiências práticas e/ou a falta de socialização no universo teórico solicitado no enunciado da questão, levou a que a maioria dos estudantes desconsiderasse as premissas analíticas chaves da teoria social crítica, quando muito, enfatizando estereótipos ideológicos do senso comum sobre as temáticas abordadas nas questões.
A percepção aqui apresentada é representativa dos resultados da avaliação de todos os corretores.
No entanto, seria importante ter mais informações a respeito da possibilidade de boicote ao exame, suas razões e procedência.
Tais informações, talvez, pudessem contribuir para uma apuração mais aproximada das reais condições dos estudantes da área de Serviço Social acerca dos conteúdos programáticos, do projeto ético-político, da formação profissional e do traçado histórico, cultural, político e social do Brasil e do mundo. 12 Lei nº 10.861, de 14 de abril de 2004.

% matched lemmas: 
\end{textsample}
